\chapter{一废太子}
% 原书第31页

说到康熙废太子，伴随而来的第一个问题就是为什么废黜太子？从太子1675年册立算起，到1708年被废时，已历三十三年。到底发生了什么事情，让康熙做出了这个决定？现在要了解这件当时“变天”的大事，可以看到的资料并不多。尽管废黜太子是当时昭告天下之事，但在文人笔记或野史资料中鲜有记载；或者也许记载过，但后来又因为各种原因没有留存下来。对于废黜太子期间发生的事，最主要的来源就是康熙朝《实录》对废黜太子当天康熙各种言辞约一千一百字的记载（见附录一）。


% 原书第31–32页

在这份关于废太子事的官方记录里，废黜太子的原因分成了两类：老毛病和新问题。前三段痛陈太子的老毛病，总结起来就是太子骄奢淫逸、积习难改；第四段从太子对皇十八子漠不关心开始，讲太子最近发生的新问题。从论述的长度上看，太子的各种老毛病是这次废黜太子的主因，也就是这份记录最开头写的：“朕包容二十年矣，乃其恶愈张。”后来的学者在叙述废黜太子之时，都沿着这段记录的叙事框架，把废黜太子的原因解释为一个日积月累长期矛盾的结果。总结起来就是康熙此前已对太子有诸多不满，而且还感到了太子对继位的渴望，让他觉察到了威胁；近期一些蛛丝马迹又进一步触发了他的警觉，因而决定废黜太子。\NoteRef{02}{1}


% 原书第32页

1970年代，有学者还在康熙身边著名汉臣李光地之孙编辑的《李文贞公年谱》中发现，李光地提到的“帐殿夜警”，是太子被废后官场上流传的废黜太子的导火索。\NoteRef{02}{2}“帐殿夜警”与《实录》中“伊每夜逼近布城，裂缝向内窃视”的记载相呼应，意思是康熙感觉太子晚上从其营帐外的缝隙中窥探自己，让他感受到了威胁。也就是说康熙认为太子要暗害他，这才下了决心要废黜他。李光地的说法相当于为《实录》中的记载划了重点：康熙与太子的关系已经疏远到了康熙怀疑其暗害自己的程度。后来无论是学者的研究，还是小说、影视的改编，都按照这个思路来解读，认为康熙与太子之间关系出现问题由来已久，而康熙废太子是他们关系长期恶化后的一个结果。


% 原书第32–33页

如果说康熙废太子是冰冻三尺，非一日之寒，那为什么康熙很快又反悔，两个月后就开始张罗复立太子？现在基于《实录》建构出的历史叙述其实无法解释这个问题，太子真有那么多新老问题，康熙又怎么会那么快改变决定？而且哪怕在群臣反对下，也要力排众议，坚持复立太子。


% 原书第33页

1709年4月18日康熙复立太子时，《实录》也给出了康熙的解释。不过这次，废黜太子的原因变成了废太子之前得了“暴戾狂易”的病：


% 原书第33页

\begin{bookquote}（康熙四十八年三月初九日庚辰，1709年4月18日）以复立皇太子允礽，……建立嫡子允礽为皇太子历有三十余载，不意忽染暴戾狂易之疾。臣以祖宗所遗洪业，及万邦民生，不敢姑息，遂行退废，绝无纤毫私意，前者告祭俱已悉陈。\end{bookquote}

% 原书第33页

在复立的诏告中，逻辑变得十分直接简单。废太子之前得了病，现在病治疗好了，所以就恢复了太子的位置。对于太子得病，《实录》中有简单提及。这个病是巫术所致。


% 原书第33–34页

换句话说，《实录》在太子被废以及复立之时，两次给出的解释是不一样的。在宣布复立太子时，废黜太子的理由变成了一个临时性的。以前的学者都注意到了前后说法的不一致。为了弥合分歧，学者们替《实录》做了一个推测性的解释，就是康熙废太子后，后悔了，认为太子位空缺导致的皇子夺嫡问题更麻烦，因而借着太子得病这件事找了个台阶来复立太子。也就是在现在的历史认识中，太子得病只是康熙复立太子时的一个借口。\NoteRef{02}{3}


% 原书第34页

尽管康熙朝储位之争是学界内外都极为关心的一段历史，但限于资料，现在对第一次废黜太子的认识与民国时期相比其实并没有大的改变。幸而新近出版的传教士纪理安的《北京纪事》为我们重新认识这段历史打开了一扇独特的窗口，其叙述逻辑总结起来有两点：一是废黜太子是偶然事件；二是康熙决定复立太子的转折点就是太子的病被治好，并不是借口。


% 原书第34页

纪理安用很肯定的口气写下了他的两点判断。他能有这样自信又源于两点：一是对传教士来说，皇帝也好，太子也罢，都不是历史人物，而是身边的人。他们在长期接触中，了解康熙与这些皇子的关系。二是，在康熙决定废太子的那次行程中，传教士巴多明和一位西医医生就在行营中。就在废太子事发生的前几天，他们还在和康熙以及太子打交道。纪理安为此事专门做过调查，并让巴多明写下了详细的记录。也就是说，对康熙与皇子的长期了解，加上对事件发生时候情况的具体分析，让他对废太子事件做出了自己的判断。


% 原书第34–35页

在《北京纪事》10月15日条目下，纪理安记录了皇十八子去世的消息，并说康熙当时还没有宣布死讯。这条消息是他几天后加到15日下面的。他在10月16日收到了跟在康熙身边的巴多明神父前一天从行营中发来的信件。在这封信中，除了传递皇十八子去世的消息以外，纪理安称巴多明还写了一件“不可思议”的事：


% 原书第35页

\begin{bookquote}还发生了一件可怕的悲剧，在欧洲的历史编年史中都很难看到相似的事件。不过我现在还是不谈论。等到我直接了解过后，11或者12月再来写这件事。\NoteRef{02}{4}\end{bookquote}

% 原书第35页

在所有的清史记录中，皇十八子因病去世和皇太子被废都是同一天宣布的。所以纪理安这里说的悲剧，就是废太子事件。从他的用词上可以看出他收到巴多明神父来信时的惊讶和难以置信，因而他需要等巴多明回来当面确认后再写到记录中。在《实录》里，废黜太子和皇十八子病逝写在10月17日，和传教士的记录很接近。


% 原书第35页

纪理安经多方问询后，在《北京纪事》1708年10月27日的记录中写下了他了解到的废太子事件，并专门做了评论，认为太子是被陷害的：


% 原书第35–36页

\begin{bookquote}对太子的主要指控绝对是虚构的，尽管现在很多人相信了，但那是因为很多人重复说到了这个指控。大皇子一党把场景编造得很好，就算谨慎的人也很难不跳进他们的故事中。流行的说法是太子在晚上拿着匕首进入到他父亲的帐篷中，意图行刺。流言精心编造了帐篷周围的情景，就连我们一些跟去的神父都觉得不得不信。\NoteRef{02}{5}\end{bookquote}

% 原书第36页

纪理安记录跟李光地的说法一致，也把废黜太子的导火索放在了夜间行刺上，同时也证明了，确实如李光地所言，这个理由是当时广为流传的说法。但同时纪理安用很肯定的口气断言太子是被大皇子一系陷害的。纪理安能有如此肯定的判断源于神父巴多明的观察。废黜太子期间，巴多明就跟在康熙与废太子的行营中。他的记录为我们提供了来自现场的信息。

% 读者评论：02 第二章 一废太子.md，第1条
\begin{comment}
巴多明当时在康熙行营中，担任翻译。他参与皇十八子的诊疗交涉，多次与康熙、太子、大阿哥面对面谈病情，随后继续留在行营。

何多敏当时与巴多明同在行营，是负责治疗皇十八子的外科医生。亲自诊察、手术，也与康熙和皇子交涉；停止治疗后仍留在行营。

纪理安当时留在北京，没有随这次行营在现场。他收到了巴多明来信，等他们回来后询问经过，再结合宫中消息写成记录。
\end{comment}


\begin{sectionnotes}

% 原书注释：原书第392页
\BookNote{02}{1}{王锺翰，《清世宗夺嫡考实》，1064—1066；Wu, Passage to Power, 118-119；杨珍，《清朝皇位继承制度》，169—205。}

% 原书注释：原书第392页
\BookNote{02}{2}{Wu, Passage to Power.}

% 原书注释：原书第392页
\BookNote{02}{3}{比如孟森，《明清史讲义》，462—463；王锺翰，《清世宗夺嫡考实》，1065；杨珍，《清朝皇位继承制度》，207—209；姚念慈，《康熙盛世与帝王心术：评“自古得天下之正莫如我朝”》，130—132。}

% 原书注释：原书第392页
\BookNote{02}{4}{AP3, 458.}

% 原书注释：原书第393页
\BookNote{02}{5}{AP3,466-467.}

\end{sectionnotes}
\section{神父巴多明的报告}

% 原书第36–37页

法国耶稣会传教士巴多明（Dominique Parrenin，1665—1741）是传教士中公认的语言天才。他1698年入华，很快就把汉语和满语的听说读写全部学会了。1706年，教皇使节多罗到北京时，他是康熙身边主要的翻译。语言上的天赋使他在京城的权贵圈以及内务府中建立了各种人脉关系。雍正继位初年的敏感时期，他还能通过内务府的旧相识打探到雍正的动向，同时还能在雍正年间请到皇十三子为传教士说话。\NoteRef{02}{6}他在雍正后期和乾隆朝，还在宫中开设拉丁语课程，为清廷培养拉丁语翻译人才。\NoteRef{02}{7}康熙废黜太子期间，他正是作为翻译跟在康熙的身边。


% 原书第37页

巴多明并不是康熙离开北京去行猎时就跟出去的。《实录》中记录这次外出始于四十七年七月十八日（1708年9月2日），太子以外，尚有大皇子和皇十三、十四、十五、十六、十七、十八子随行。在途中，皇十八子生病。康熙命令外科医生何多敏（Giandomenico Paramino，1661—1713，此人不是神父）前来治疗，巴多明作为翻译随行（参看书签一时间表）。他们是1708年9月27日晚7点接到命令后出发的。\NoteRef{02}{8}


% 原书第37页

巴多明说，赶到之后，在帐篷内看到了太监抱着的皇十八子。当时肉眼所见的症状是膝盖红肿，同时从太监口中得知，中医认为已经没有内部的病了。这也解释了为什么要匆忙中把京城的外科医生叫去。由于何多敏和巴多明是被召去治病的，所以巴多明报告的撰写围绕着治病展开。从报告不厌其烦地详述各种细枝末节来看，当时北京耶稣会传教士面临很大的压力，毕竟他们参与医治的皇子最后去世了。当时北京还有一位西洋的内科医生，他已经向当时住在山东的北京教区主教汇报过这次医疗事故，同时声称这次医疗事故让康熙疏远了传教士。这位西洋内科医生是教皇特使多罗带到北京的，因而派系上讲，与众北京耶稣会士有矛盾。\NoteRef{02}{9}


% 原书第37–38页

第二天早上，何多敏和巴多明再次去给皇十八子看病。这次是在白天的光线下。何多敏仔细给皇十八子做了检查后，发现他十分虚弱，发着烧，舌苔干燥，同时除了膝盖（knee）有肿胀外，下颚（jaw）位置也有问题。何多敏告诉当天负责皇十八子治疗事宜的大皇子，他认为皇十八子已经无法医治了。巴多明说大皇子和当天听到这个意见的人都不敢告诉康熙。


% 原书第38页

9月30日上午，康熙把二人找到面前，询问何多敏，鉴于皇十八子膝盖肿的部分已呈黑色，能不能用刀割掉或者用针刺破？何多敏说皇十八子身体虚弱，禁不起手术，而且他认为手术后病情也不会有好转。康熙让他们离开了。下午时，太子和大皇子一起来找他们，再次询问他们能不能把皇十八子膝盖的肿胀部分切除。何多敏医生还是重复说明切除对病情没有帮助；于是两位皇子传达了康熙的意见，说中医认为切除后就能有好转；康熙命令何多敏手术，并告诉何多敏不要有心理负担，决定是他做的。何多敏给皇十八子做了切除手术，并用上了从埃及带来的一种药来缓解手术时的疼痛。\NoteRef{02}{10}


% 原书第38页

八月二十日（10月3日，《北京纪事》采用阳历，但巴多明记录用的是阴历），他们与太子、大皇子和康熙一起讨论了病情，结论是皇十八子的病情没有好转。八月二十一日，他们再次被召去诊病，这时何多敏告知太子，他认为皇十八子的病已经没有希望了。当时距离皇十八子真正病逝还有半个月时间。\NoteRef{02}{11}


% 原书第38–39页

何多敏没有再被召去治病，但他和巴多明就留在行营中，由此成为了观察者。巴多明对后来皇十八子的治疗和病情的描述与零星留存的满文奏折正好可以互相印证。在何多敏做手术的同时，康熙在八月十九日写了满文谕旨传回北京，晓谕皇三子和皇四子，让他们去找“刘妈妈”和“外科大夫妈妈何氏”，并把她们送来给皇十八子治病，并嘱咐皇三子办事时不要“张扬”。\NoteRef{02}{12}皇三子在回奏中还专门写到他们是八月二十日未时头三刻（相当于现在下午1：30—1：45）收到奏折，以此凸显他们是在争分夺秒地遵照吩咐在找人。他们找的人在八月二十二日送到了行营。巴多明记录中说，在何多敏说不能医治后，有女医生从北京赶来给皇十八子治病，这和满文档案的记载吻合。\NoteRef{02}{13}而且在八月二十四日一度传言治疗有效果。巴多明和何多敏被太监叫去看望好转的皇十八子。巴多明说这是太监要羞辱他们，意思是他们的判断错了。巴多明提到病情一度好转，又迅速恶化，这一点也与现存的满文记录吻合。皇三子于八月二十九日收到康熙谕旨，其中提到皇十八子病情“渐次好些”。但到了九月初二，康熙又写“初二日晨突变，病势甚重”\NoteRef{02}{14}。


% 原书第39–40页

除了从北京新来的女医生外，巴多明称喇嘛也在参与治疗。喇嘛做法以后，要求全体人员不得骑马，说否则对皇十八子病情不利。于是太子步行走在大部队最前方，带领整个队伍徒步前进。\NoteRef{02}{15}巴多明记录中提到这一天是八月二十四日，可见这个时候太子还是太子，还在作为表率带领众人前行。太子是九月初四（10月17日）宣布被废。也就是说，引起废黜太子的事件发生在八月二十四到九月初四这十天之中。巴多明的记录中没有提供废黜太子的原因分析。他只是说九月初四的中午首先宣布了皇十八子的死讯，下午宣布了废黜太子，太子被收监，并被上了枷锁。\NoteRef{02}{16}


% 原书第40页

巴多明和太子、大皇子以及康熙都有多次面对面关于皇十八子病情的探讨，而且他吃住都在行营中，在整个过程中都没有察觉到异样。所以在他眼中，废黜太子就是突然或者说临时发生的，并不是一件有预先谋划或者安排的事。康熙这次带出去的皇十七子时年十二岁，皇十八子仅八岁。而且根据巴多明的记录，皇十八子的母亲也在这次行程中。这些安排说明这次出行就是单纯的打猎。在出发前，康熙并没有提前考虑过要废黜太子。太子在行程中的大部分时候，还担负着各种大小事务，同样说明康熙废太子并非提前安排。


% 原书第40–41页

留在京城的纪理安用十分肯定的语气判断太子是被大皇子陷害的。这一判断是他在向巴多明和何多敏了解了情况，并向康熙身边人打听后，外加他自己对太子和大皇子性格品行长期了解的基础上综合而得。退一步说，就算我们怀疑纪理安的判断，但最起码这个判断表明了纪理安极度不相信太子会行刺康熙，因而才说太子是被陷害的。首先，纪理安和巴多明与太子相熟。比如在两年前，纪理安就是因为有两周时间天天一大早就到太子府上去，待到日落才回教堂，因而受到当时在京教皇使节的质问，认为他不务正业。\NoteRef{02}{17}再比如废黜太子前半年，俄罗斯派了一位东正教神父到北京。太子就找到纪理安和巴多明，让他俩去跟这位新来的神父聊一聊。太子让他们去，目的是去试探这个俄罗斯人是不是真的神父，也就是说去考察一下这位新来之人的背景。\NoteRef{02}{18}这样与太子往来的小事在传教士的记录中还有很多。传教士对于太子这个人有直接了解。这是他们认为太子不会行刺康熙的基本出发点。


% 原书第41–42页

其次，他们对太子还有许多间接的了解和观察。比如废黜太子前四个月，在《北京纪事》的6月15日下，纪理安记录了一件事：杭州送来一份加急密折，内容是关于打着朱三太子旗号反清复明之事。送密折的人由于过度劳累，当场晕倒在康熙面前。康熙只让太子看了这封加急密折，并当着太子面写了回复。然后康熙让两位满人和一位蒙古人带着这封回复加急传往杭州。纪理安写下这个细节并解释说，派三位一起出发是考虑到如果某一位在途中晕倒或者体力不支，那么其他人可以继续递送。\NoteRef{02}{19}本书前面已经提到过，纪理安等人在宫中的工作地点就在养心殿附近，因而平时与内廷中的太监、奴才多有接触，经常一起说话聊天，这些人是传教士获取信息的重要渠道。像上面所说送信的人当场晕倒这样的戏剧性情节，大概率就是从宫里的太监处听来的。什么事紧急到传信的人能当场晕倒，宫里的人都会好奇。传教士也不例外，同样对八卦有兴趣。记录这则消息时，太子还没有被废。在传教士的眼中，康熙极其重视太子，这样最私密最重要的信息没有知会任何官员，而只给太子看了。这些细节对于传教士来说，是康熙倚重和信任太子的明证。有了这些大小事的背景，就可以理解为什么纪理安会在记录中以很肯定的口气判断太子是被陷害的，并能准确指出加害人就是大皇子。


% 原书第42页

大皇子也是传教士熟悉且打过很多交道的皇子。在废黜太子之前的几年，康熙派大皇子管理内务府中跟传教士相关的事务，因而传教士和大皇子有很多接触。1706年教皇使节入京后，大皇子也是经常出现在《北京纪事》中的人物。在传教士笔下，大皇子是一位性格直率，或者说性格粗鲁的皇子。他可以当着众人的面就开始发火，甚至在纪理安的记录中，专门用文学的方式描绘过他说话的场景。比如康熙曾经派白晋神父给教皇送礼物，而教皇使节也派了一位叫沙国安的神父加入白晋使团中。二人到了广州后，开始为谁是使团领导起了争执，并把争执传回了北京。大皇子听说后，当着很多神父就发起火来：


% 原书第42–43页

\begin{bookquote}这怎么可能是个问题？怎么可能有人质疑白晋不是我们的使节？他在我们宫中二十年，满文汉文都会，一直在父皇那里效力。而谁是沙国安？他谁啊？谁认识？\NoteRef{02}{20}\end{bookquote}

% 原书第43页

就在太子被废前几个月，地方上的传教士还因为内务府发给他们的身份证明遇到了麻烦。地方官认为传教士作为僧侣应该归礼部管，因而当他们看到传教士出示的由内务府签发的身份证明，就怀疑是伪造的，把传教士告到了礼部。那时候负责协调礼部和内务府的就是大皇子。像前面还讲到过，大皇子和徐日昇神父之间，是前者可以拿后者的胡子开玩笑的长期亲近关系。简单来说，当时的传教士和太子、大皇子都有数十年的往来接触。当他们在评点康熙废太子时，跟后世史官根据各种记录评判一桩政治事件不同，更像是一位家里人在谈论大家庭中的家长里短。正如《红楼梦》中绘声绘色写到的焦大对贾府的点评：


% 原书第43页

\begin{bookquote}我要往祠堂里哭太爷去。那里承望到如今生下这些畜生来！每日家偷狗戏鸡，爬灰的爬灰，养小叔子的养小叔子，我什么不知道？\end{bookquote}

% 原书第43页

焦大“哭太爷”的本质就是资历的展示——他是家里人，因而他“骂”出来的事，就是他作为家中的亲历者平日看到、感受到的。而传教士从身份上说也是康熙府中之人，跟康熙家中上下都有联系。纪理安对于太子行刺的评论就是一位熟悉康熙家事的人做出的评论，所以他的论断十分肯定。


% 原书第44页

读过纪理安的记录以后，再重读李光地的记录，就会发现李光地这样的康熙近臣在听到太子行刺后的反应，其实本质上与纪理安是一样的。《李文贞公年谱》中的记录是：


% 原书第44页

\begin{bookquote}还宫有顷，召公独对。公顿首请曰：“帐殿之警，上果稔知其主名？必无刺谬乎？”上默不答，怅然久之，曰：“此直为鬼物所戏耳，何丧心至是？”\NoteRef{02}{21}\end{bookquote}

% 原书第44页

康熙一回到北京就召李光地单独见面。见面前，李已经听到过行刺的传言，因而当他问询康熙这件事的时候，是经过思考的。这个问题本身就已经显示出李光地不相信太子行刺，因而他用了两个反问句来质疑康熙的判断：帐殿之警事主是谁，真的确定吗？一定不会错吗？李光地这样的汉臣对皇上说话向来恪守上下尊卑，极其克制，这样连用两个反问来质疑康熙的判断，已经很明确地表达了他的怀疑。康熙也确实拿不准，所以才会“默不答，怅然久之”。


% 原书第44–45页

纪理安和李光地的第一反应都是觉得太子行刺不可信。最起码，他们的怀疑说明他们之前没有觉得康熙和太子的关系有大的问题；或者反过来说，如果康熙和太子之间以前就有大的问题，积累已久，那么纪理安和李光地也不会有这么大的反应。特别是李光地，他应对康熙，从来都是极其小心。如果他曾察觉到康熙与太子之间关系早就有大问题，就不会这样反问。就像一对夫妻突然宣布离婚了，谁会惊讶呢？一直认为他们关系没什么问题的人才会惊讶。


% 原书第45页

在11月2日和3日的条目下，纪理安提到他看过当时正在官场流传的关于废黜太子问题的诏书。《实录》中也确实提到康熙在11月回京后，又几次在大臣面前谈到过废黜太子之事，而且还到太庙等地昭告天地。同时11月初（九月二十四日），还公开发布了废黜太子的诏书，内容是在前面提到康熙行营中所说版本的基础上简化而成。从程度上看，诏书里提到的问题比行营中提到的问题性质要轻许多。而且这时给出的理由已经倾向于太子生病了。\NoteRef{02}{22}纪理安对诏书的评价是，其中列举太子的问题是康熙为了要向天下证明（justify）他废太子是有道理的：


% 原书第45页

\begin{bookquote}就算这些（公开太子的种种问题）能够让众人都对太子有看法，但这些也同样可以用来批评皇帝，以及让整个皇家蒙羞，因为他们没有把太子教育好。而且，作为比较，事实上太子的兄弟们在这些提到的方面上，并没有做得更好。\NoteRef{02}{23}\end{bookquote}

% 原书第45–46页

纪理安没有具体写下到底他当时看到的版本中公开列举出的理由有哪些，但整体上他并不认为诏书中太子被废的理由有说服力，因为诏书中提到的问题在他看来是当时皇子们普遍存在的毛病。纪理安还提到一个问题：康熙废掉的太子，也是他一手培养的儿子，他公开把儿子说得越差，问题越大，其实对康熙来说也是自取其辱。“子不教，父之过”，纪理安点出了这一层意思。写完这一段评价后，他把废太子事件发生的矛头指向了大皇子。


\begin{sectionnotes}

% 原书注释：原书第393页
\BookNote{02}{6}{《耶稣会士中国书简集》，卷2，327。}

% 原书注释：原书第393页
\BookNote{02}{7}{学习拉丁文见郭福祥，《巴多明与清官西洋学馆》。}

% 原书注释：原书第393页
\BookNote{02}{8}{AP3, 472.}

% 原书注释：原书第393页
\BookNote{02}{9}{AP3,467.}

% 原书注释：原书第393页
\BookNote{02}{10}{AP3,476-477.}

% 原书注释：原书第393页
\BookNote{02}{11}{AP3, 477.}

% 原书注释：原书第393页
\BookNote{02}{12}{《康熙朝满文朱批奏折全译》，597。}

% 原书注释：原书第393页
\BookNote{02}{13}{AP3,478.}

% 原书注释：原书第393页
\BookNote{02}{14}{《康熙朝满文朱批奏折全译》，600。《实录》这一天（乙亥日）也是如此记载：“自十八阿哥患病以来，朕冀其痊愈，昼夜疗治。今又变症。”}

% 原书注释：原书第393页
\BookNote{02}{15}{AP3,478.}

% 原书注释：原书第393页
\BookNote{02}{16}{AP3, 478.}

% 原书注释：原书第393页
\BookNote{02}{17}{AP1, 363.}

% 原书注释：原书第393页
\BookNote{02}{18}{AP3,42–43.}

% 原书注释：原书第393页
\BookNote{02}{19}{AP3,204-205.}

% 原书注释：原书第393页
\BookNote{02}{20}{AP1,378.}

% 原书注释：原书第393页
\BookNote{02}{21}{《李文贞公年谱》，194—195。}

% 原书注释：原书第393页
\BookNote{02}{22}{《圣祖仁皇帝实录》：丁酉。以废皇太子允礽，颁诏天下：“（此处省略关于江山社稷的大论）允礽自立为皇太子，时勤教谕，并简名望大臣为之讲明性理，历有年所。顾秉性乖戾，罔体朕心，违背朕训。虽愆尤日积，尚冀其悔悟自新。屡次南巡江浙，西巡秦晋，皆命允礽随行。原望其谙习地方风俗，民间疾苦，乃辄强勒督抚大吏及所在官司，索取财贿。所用宵小匪类，尤恣意诛求，肆行攘夺。夫地方物力，皆属小民脂膏。朕屡谕允礽，宜加节俭，伊乃穷奢纵欲，逞恶不悛。既已苛索外吏，复夺取外藩入贡马匹等物。私用内外库帑，为数甚多。流毒臣民、安所底极。迩来暴虐慆淫，过端弥著。自诸王以及大臣官员，悉被非礼凌辱，横加捶挞。向因索额图、常泰，交通设谋，朕洞觉其情，置索额图于死，而允礽时蓄忿于心。近复有逼近幔城，裂缝窥伺，中怀叵测之状。凡此举动，类为鬼物所凭，狂易成疾。《书》曰：‘天视自我民视，天听自我民听。’人心所恶，天必厌之。宗社事重，何以承祧？朕图维再三，万不获已，于康熙四十七年九月十八日，奉皇太后慈命告祭天地、太庙、社稷，特废黜拘禁。……”}

% 原书注释：原书第394页
\BookNote{02}{23}{AP3, 480.}

\end{sectionnotes}
\section{大皇子背后的故事}

% 原书第46页

在纪理安的记录中，康熙一行于10月29日早上8点回到北京，废太子被关押在上驷院（养御马的地方）一侧，由大皇子与皇四子负责看押。\NoteRef{02}{24}这里记录的地点和两位负责看押的皇子都和《实录》所记一致。\NoteRef{02}{25}在11月6日的条目下，纪理安记录朝廷正式宣布“皇太子”的称号改为“二阿哥”。也是在这一天下面，纪理安说朝堂上下已经有很多人陆续忙着去接近大皇子，还有人已经在暗中恭喜大皇子。\NoteRef{02}{26}


% 原书第46–47页

大皇子不是康熙真正的长子，在他前面其实还有四位，只是都早夭了。从现有记录来看，大皇子善于骑射，早年康熙平定准噶尔的远征，他都参与了。大皇子生于康熙十一年（1672），废黜太子这年，他三十七岁，比太子大两岁，比皇三子大五岁，比皇四子大六岁，比皇八子大九岁，比皇九子大十一岁，比皇十三子大十四岁。后来在众皇子中崭露头角的皇十四子，当年二十一岁。根据传教士的记录，这一年正好康熙在安排为他在宫外修建独立的皇子府。另外，除了太子的母亲出身高贵以外，本书涉及的其余几位皇子的母亲出身都比较低微，他们命运的发展跟他们的母亲是谁并没有关系。\NoteRef{02}{27}

% 读者评论：02 第二章 一废太子.md，第2条
\begin{comment}
几位夺嫡核心皇子是这样的，这些皇子有：皇长子\textasciitilde{}皇四子、皇八子、皇十四子。康熙皇十子母家出身很好，但皇十子不是夺嫡核心。
\end{comment}


% 原书第47页

太子被废后，大皇子迎来希望，也是地方上的传教士从公开张贴的告示中得出的印象。住在江西的传教士殷弘绪（Francois Xavier d'Entrecolles）神父在一封写给教会上级的信中就说，他看到了“公开张贴的告示中充满了谴责皇太子行为的声明和责骂”，同时他也看到了“称为大阿哥的大皇子”获得了赞扬，而且他还说大皇子“以为自身地位不久就会因为兄弟垮台而提高了”。也就是说，地方上的传教士看告示后也感觉大皇子是这次废太子事件的受益者。\NoteRef{02}{28}


% 原书第47页

由于大皇子之前主管内务府有关西洋人的事务，在《北京纪事》中，他是比太子出现频率还要高的皇子。也是由于这个原因，纪理安对大皇子身边发生的事记录得也比较详细。当写到大皇子接受各方恭喜时，纪理安评价道：


% 原书第47页

\begin{bookquote}他难以抗拒那些恭喜。在这些恭喜的簇拥下，他很傻地放下了对忠于或者支持太子的人的戒心。\NoteRef{02}{29}\end{bookquote}

% 原书第48页

纪理安笔下呈现出来的大皇子的状态，用中国传统的说法，就是被胜利冲昏了头脑，有点飘飘然，忘乎所以了。纪理安说大皇子为了争夺皇位，在想办法让其他皇子在康熙那里出问题。此处，纪理安还专门点明“特别是皇八子”\NoteRef{02}{30}。这表明太子被废后，大皇子看到的最强竞争者是皇八子。这与《实录》中呈现出的大皇子支持皇八子的说法相悖。

% 读者评论：02 第二章 一废太子.md，第3条
\begin{comment}
《圣祖实录》卷233

拘禁允礽时，允禔奏：“允礽所行卑污，大失人心。相面人张明德曾相允禩，后必大贵。今欲诛允礽，不必出自皇父之手。”

紧接着，康熙的反应是：

“倘果同允禩聚集党羽，杀害允礽……岂暇计及于朕躬有碍否耶。”
\end{comment}


% 原书第48页

纪理安还记录了一件让大皇子最终失势的事件。这个事件在所有的清史记录中都没有提及。大皇子在被看好的情况下招惹到了本不应该招惹的人：康熙身边的太监。事情的起因是大皇子打死了一名废太子身边的太监，这引起了整个太监群体的公愤，于是康熙身边的太监作为太监群体领袖，用计扳倒了大皇子。据纪理安记载，康熙身边的太监并没有把大皇子打死废太子太监的事直接告知康熙，因为这样会让康熙认为他们的报告有私心，从而降低他们说话的分量。他们选择把大皇子从前打死一位宫内工匠的事报告给了康熙。在纪理安的记载中，康熙身边的太监在报告时，察言观色，根据康熙的反应还添油加醋了一番，说这位工匠之死让宫里所有工匠都深感不安。最后还说现在大臣们都惧怕大皇子。\NoteRef{02}{31}


% 原书第48–49页

康熙听完汇报后，把大皇子召到养心殿中，当面质问，最后解除了大皇子的所有职务。当天宫里还有很多人在等待召见，一些传教士也在场。纪理安还提供了细节，说大皇子从康熙那里出来时涨红了脸。当时还有朝臣在等他，准备与他谈事，他都没有搭理。他直接走到太监总管那里，在后者耳边说了几句话。太监总管听完后立马变了脸，流出很多眼泪。纪理安生动地写出了太监转瞬即来的眼泪，带有一种幽默的讽刺，因为他还特别点明他们传教士都知道，这一切本来就是这些太监们安排的。纪理安说大皇子没有听太监的安慰，也没有和其他人说话，就匆匆离开了。\NoteRef{02}{32}在纪理安的叙事中，大皇子在经过这件事后便失去了竞争储位的资格。


% 原书第49页

回到跟本书相关的问题，如果真有太监扳倒大皇子之事，为什么《实录》不记录呢？这不大可能是一个偶然的疏漏。要知道皇四子是和大皇子一起负责看守废太子的人，当时宫里发生的一切，皇四子都近在咫尺，不可能在编辑《实录》的时候，完全遗忘了这件让大皇子失势的大事件。这其中定有隐情。而这个隐情需要在了解了夺嫡斗争中的下一场，也就是和皇八子密切相关的看相人张明德案以后才会知晓。


\begin{sectionnotes}

% 原书注释：原书第394页
\BookNote{02}{24}{AP3, 470.}

% 原书注释：原书第394页
\BookNote{02}{25}{《圣祖仁皇帝实录》，卷234。}

% 原书注释：原书第394页
\BookNote{02}{26}{AP3, 482.}

% 原书注释：原书第394页
\BookNote{02}{27}{康熙家里的儿子情况，参看杨珍，《康熙皇帝一家》，第五章“皇子”。}

% 原书注释：原书第394页
\BookNote{02}{28}{《耶稣会士中国书简集》，卷2，35。}

% 原书注释：原书第394页
\BookNote{02}{29}{AP3,482.}

% 原书注释：原书第394页
\BookNote{02}{30}{AP3, 482.}

% 原书注释：原书第394页
\BookNote{02}{31}{AP3,482-483.}

% 原书注释：原书第394页
\BookNote{02}{32}{大皇子之前负责管理传教士，与传教士的接触很多，在《北京纪事》前期的记录中有很多关于他们之间打交道的记录。}

\end{sectionnotes}
\section{张明德案中被篡改的历史}

% 原书第49–50页

10月底回宫后，康熙忙于处理各种废黜太子后面临的实际问题。比如太子的称谓、住所等等都要更改。11月初，皇子间的夺位之争已经悄然拉开了序幕。争斗最早的记录是在康熙四十七年九月二十五日戊戌（1708年11月7日）。\NoteRef{02}{33}这一天，康熙指示相关官员审理张明德案。起因据《实录》记载，是大皇子在康熙面前推荐皇八子，说相面人张明德给皇八子看过相，说他“后必大贵”。\NoteRef{02}{34}


% 原书第50页

相面人张明德案是储位之争的第一案。小说、电视涉及储位之争，对这个案子都有演绎和发挥。简单来说，张明德是一位以看相来算命的术士，游走于满人皇亲国戚之中。在太子被废之前，皇八子就找张明德算过。根据后来雍正朝的庭审记录，皇八子当时没有儿子，他自陈算命主要是为了求子。\NoteRef{02}{35}就明清之际社会整体而言，看相、算卦、测八字是日常行为，并不是什么见不得人之事。算命本身不是这个案子的核心，但废黜太子事件发生后，皇八子参与算命却成了皇八子的反对者搞垮他的事由，成为了朝堂上的案件。


% 原书第50页

就现有的历史记录来说，无论是正史还是后来的文学发挥，这个案子的主角都是两位皇子：大皇子和皇八子。整个案子的审理分为两个阶段。第一阶段：康熙认为相面人张明德算命惑众，因而需要处理他。康熙给官员交代：


% 原书第50页

\begin{bookquote}闻彼曾为允禩看相，又散帖招聚人众，其情节朕知之甚明。此案甚大，干连多人，尔等慎毋滋蔓，但坐张明德一人审结可也。\NoteRef{02}{36}\end{bookquote}

% 原书第51页

康熙给出的指示是“慎毋滋蔓”，把张明德一人问罪就行了。他知道张明德是给皇亲国戚算命的，所以特别叮嘱别把事情复杂化了，以至于“干连多人”。在审问张明德案之始，皇八子并不是审讯的目标。《实录》中还记载，皇八子这时正被康熙派去审问原任内务府总管的家产案。在11月10日（九月二十八日辛丑），他还在向康熙汇报案情。也就是说皇八子此时尚未涉案。


% 原书第51页

张明德案的审问在11月11日这一天变了调。张明德被查出之前有谋反刺杀太子的计划（该计划是在太子被废之前）。更重要的是，张明德曾把计划讲给过皇八子，但皇八子并没有采取任何措施阻止，也没有上报过康熙。还有其他的满人权贵也听闻过张明德的计划，也都没有向康熙汇报过。这样一来，张明德的案子就有了谋反性质。大学士温达等报告的案情记录在《实录》九月二十九日壬寅（11月11日）下面，他们建议对张明德斩立决。同时这一天下面还写到，康熙为此事“召诸皇子入乾清宫”。康熙当着众皇子面，责骂了皇八子：


% 原书第51–52页

\begin{bookquote}着将允禩（皇八子）锁拿，交与议政处审理。皇九子允禟语皇十四子允禵云：“尔我此时不言何待。”允禵奏云：“八阿哥无此心，臣等愿保之。”上震怒，出所佩刀，欲诛允禵。皇五子允祺，跪抱劝止。诸皇子叩首恳求。上怒少解，命诸皇子挞允禵，将允禟、允禵逐出。\NoteRef{02}{37}\end{bookquote}

% 原书第52页

这里还涉及皇九子和皇十四子，并由此成为他们二人与皇八子交好的证据。\NoteRef{02}{38}在皇八子被“锁拿”这一天后，案情就进入了第二阶段：皇八子是否参与谋害太子成为审讯的重点。


% 原书第52页

《实录》四十七年十月初二日甲辰（1708年11月13日）下记录了新一轮调查结果，确认了张明德确实向皇八子提到过刺杀太子之事，以及其他几位权贵也找张相过面，也听他说过刺杀太子的谋划。最后，康熙宣布将张明德凌迟处死，而对皇八子的处理是：


% 原书第52页

\begin{bookquote}贝勒允禩，闻张明德如许妄言，竟不奏闻。允禩革去贝勒，为闲散宗室。\NoteRef{02}{39}\end{bookquote}

% 原书第52页

张明德案本身与康熙废太子毫无关联，因为张案发生在废黜太子之前。张明德除了说过行刺太子以外，没采取过任何具体行动。皇八子在被审问时的解释是，当他听到张明德的谋划后，认为他是疯子，就再没有和他往来，当然他也没有进一步报告。所以，这件事突然出现在朝堂上，是被人有意翻出来的往事，其目的很明显，就是要扳倒当时皇亲国戚以及大臣们最看好的皇八子。


% 原书第53页

后来研究这段历史的学者，都明白这个案子的动机，也知道其本质是针对皇八子的夺嫡之争。作为夺嫡的第一案，这里涉及到了各皇子最初的反应，是判断各位皇子派系的依据。而且该案算是《实录》里记录详细的案子，事件的发展逻辑和各皇子的人设都有所呈现。历史上没有学者质疑过该案。由于《实录》中还载有对话，以及康熙气到拔刀砍皇十四子、皇五子抱住康熙等戏剧性场景的描写，这个桥段还经常出现在各种文学化改编中。但就算加上了各种想象和创作，也都没有人想到过这个案子的主角有可能是皇四子，也就是后来的雍正皇帝。


% 原书第53页

纪理安分别在11月6日和11月12日两个日期下记录了这个算命师的案件（传教士记录中没有提到算命师的名字）。在11月6日下面，纪理安写道：


% 原书第53–54页

\begin{bookquote}在这场闹剧中，更让人伤悲的事发生了，对皇室构成了真正的威胁。正如我所说，在所有兄弟中，大皇子最害怕的便是皇八子。众人普遍认为，皇八子是上天赐予治理国家的英才。为了除掉皇八子，大皇子故意激起皇四子的嫉妒之心，唤醒他内心深处对权力的渴望。为实施诡计，皇四子找到一名算命人，该人预言第八子将来会成为皇帝。第八子没有拒绝这一预言。于是皇四子便联合其他兄弟（纪理安注：大皇子故意回避了，没有亲自参与）指控第八子怀有篡位之心，意图谋反。\NoteRef{02}{40}\end{bookquote}

% 原书第54页

这是纪理安记录的该案缘起。他明确讲出这个案子是皇四子告到康熙那里去的，而且还特地标注了大皇子有意回避参与其中。在纪理安的叙事逻辑中，这时候大皇子回避参与也是情理之中的事，因为他刚被太监扳倒，受到康熙处分，这时完全不适合出面。


% 原书第54页

纪理安11月12日下再次记录这件事，内容对应《实录》11月11日康熙召见众皇子的记录。纪理安说：


% 原书第54页

\begin{bookquote}审讯很明确地表示，算命师是由皇四子，而不是皇八子带来的。因而皇八子的问题就变得没有那么严重，同时皇四子被他的父亲打了一顿。\NoteRef{02}{41}\end{bookquote}

% 原书第54页

纪理安的这条记录和《实录》尽管都在说张明德案的第二阶段，但内容完全不一样。皇四子变成了这一天康熙召见惩办皇子的核心人物。在《实录》中被康熙打了的人是皇十四子，而纪理安记录中被打的是皇四子。


% 原书第54页

《北京纪事》11月14日记载了皇八子被释放，同时被削去了爵位。这一条对应《实录》中11月13日的记录。


% 原书第55页

皇四子在纪理安的记录中是整个储位之争始作俑者。不仅他被康熙打了一顿，整个事件的人物关系也发生了改变。皇八子成了最大的受害者。他什么都没做，却因为皇四子的争位之心而被翻了旧账，且因此获罪。在这个案子中，大皇子和皇四子才是互相勾结。而这一点似乎也最合理，因为他二人正好被安排在一起看守废太子。


% 原书第55页

关于皇四子被康熙打了的事，《北京纪事》并非孤证。在纪理安公证的一封由江西神父在11月26日写给在澳门的教皇特使的信中，在汇报教会事务的同时，这位神父还顺带提到了北京的混乱局面。除了写太子被抓以外，还写到了“皇四子被鞭打”。也就是说皇四子被打在当时并不是秘闻。\NoteRef{02}{42}


% 原书第55页

根据学者杨珍考证，在废黜太子以前，皇四子与皇八子的关系不错。两人的府邸相连，现在的雍和宫就包括了皇八子的府邸。\NoteRef{02}{43}纪理安有一句评语，说二人因为这件事，从“友谊”变成了“敌对”，这也印证了杨珍的说法。但无奈权力的诱惑实在太大，皇四子举报皇八子之时应该也没有料到，结局是他和皇八子两败俱伤。\NoteRef{02}{44}

% 读者评论：02 第二章 一废太子.md，第4条
\begin{comment}
（一）「江西神父在11月26日写给在澳门的教皇特使的信」

出自 Acta Pekinensia 第三卷，章“November 1708”，节 21\textasciitilde{}23。

（二）「纪理安有一句评语，说二人因为这件事，从“友谊”变成了“敌对”」

出自相同文献第 6 节。

（三）捋一下老四、老八关系转冷的时间线

老四跟老八、老九、老十比邻居住，见于《康熙朝满文朱批奏折全译》第 1078 件，“胤祉奏请指定建房地折”，日期：康熙四十六年三月二十日。

b站 “王天目” 的视频 BV1aLeb6sEZM 提到老四曾经是八爷党，程度待查，详见 ext/秦道然供单与一废太子转奏考据.md 与 ext/上谕内阁与上谕八旗中的皇子交往考据.md
\end{comment}


% 原书第55–56页

明白了皇四子是举报人后，也就可以回答上一节留下的问题：为什么《实录》中没有大皇子被太监扳倒这一段？因为在《实录》的叙事安排中，大皇子需要用来顶替皇四子，充当张明德案的告发人。如果他刚刚被太监扳倒，就不适合担任这个角色。历史环环相扣，任何的改动，都会前后影响，需要做相应的改动。《实录》要把皇四子在这段历史中的印记抹去，那么就得派出另外一个人来顶替，而顶替者的相关历史也要做出相应的改动。


% 原书第56页

从历史编撰角度看，皇四子挨打这件事被雍正朝的史官删除，是中国写史传统中“笔削”手法中的“削”。但把夺嫡案的始作俑者张冠李戴，硬把皇四子改成大皇子，就是公然篡改历史，是一种虚构了。有了纪理安的记录，再读《实录》，就会发现其中有许多难以自圆其说的破绽。为什么11月13日才被削去贝勒头衔的皇八子，在两个月后就给恢复了？《实录》四十七年十一月二十八日庚子（1709年1月8日）下，记录了“复封允禩为多罗贝勒”\NoteRef{02}{45}。在现有《实录》的逻辑中，这一点是无法解释的，但却很符合纪理安的叙事逻辑：当康熙得知是皇四子把张明德引入皇家圈子的时候，就减轻了对皇八子的处罚。同时，在张明德案后，朝臣在推举新太子时，一致推荐皇八子，也从侧面说明刚刚发生的张明德案其实没有在朝堂上影响到皇八子。另外，皇四子确实对算命的兴趣远远超过一般人，他自己就回忆过：“向在藩邸，甚喜看命消遣。”\NoteRef{02}{46}他当上皇帝后，更是频繁找大臣要八字来给他们算命，并以此作为任用官员的根据。所以说，算命师张明德是他带入到皇亲的圈子中，这完全符合皇四子的特点。


\begin{sectionnotes}

% 原书注释：原书第395页
\BookNote{02}{33}{《圣祖仁皇帝实录》：“戊戌。上谕诸皇子曰：‘拘禁允礽时，允禔（大皇子）奏：“允礽所行卑污，大失人心。相面人张明德曾相允禩（皇八子），后必大贵。今欲诛允礽，不必出自皇父之手。”言至此，朕为之惊异。朕思允禔为人凶顽愚昧，不知义理，倘果同允禩聚集党羽，杀害允礽，其时但知逞其凶恶，岂暇计及于朕躬有碍否耶？似此不谙君臣大义、不念父子至情之人，洵为乱臣贼子，天理国法皆所不容者也。’”}

% 原书注释：原书第395页
\BookNote{02}{34}{《圣祖仁皇帝实录》，卷234。}

% 原书注释：原书第395页
\BookNote{02}{35}{《允禩允禟案续》，《〈文献丛编〉全编》，第3册，138。}

% 原书注释：原书第395页
\BookNote{02}{36}{《圣祖仁皇帝实录》，卷234。}

% 原书注释：原书第395页
\BookNote{02}{37}{《圣祖仁皇帝实录》，卷234。}

% 原书注释：原书第395页
\BookNote{02}{38}{据后来的一份口供显示，抱住康熙的是皇九子，而不是皇五子。当天皇九子脸被打到红肿。康熙说他们几个皇子学的是梁山好汉的义气。《允禩允禟案续》，《〈文献丛编〉全编》，第3册，137—138。}

% 原书注释：原书第395页
\BookNote{02}{39}{《圣祖仁皇帝实录》，卷235。}

% 原书注释：原书第395页
\BookNote{02}{40}{AP3, 484.}

% 原书注释：原书第395页
\BookNote{02}{41}{AP3, 484.}

% 原书注释：原书第395页
\BookNote{02}{42}{AP3, 500.}

% 原书注释：原书第395页
\BookNote{02}{43}{杨珍，《清朝皇位继承制度》，221—222。}

% 原书注释：原书第395页
\BookNote{02}{44}{AP3, 485.}

% 原书注释：原书第395页
\BookNote{02}{45}{《圣祖仁皇帝实录》，卷235。}

% 原书注释：原书第396页
\BookNote{02}{46}{《奏请赵显忠补靖远协副将折》，雍正六年五月二十一日，故宫003252号，朱批部分，台北故宫博物院清代档案检索系统。}

\end{sectionnotes}
\section{从废到立的转折}

% 原书第57页

本书序言的开头提到过，在《北京纪事》的叙述里，整个废太子案的转折点是12月4日康熙与废太子在养心殿的会面。这一天距离张明德案案发二十多天。《实录》对他们的会面也有记载：


% 原书第57页

\begin{bookquote}“……今日回宫，已召见八阿哥，并将召废皇太子一见。”顷之，内侍传谕曰：“朕适召废皇太子，亦既见之矣。自此以后，不复再提往事。废皇太子现在安养咸安宫中。朕念之，复可召见，胸中亦不更有郁结矣。”\NoteRef{02}{47}\end{bookquote}

% 原书第57页

这条记载中也说到，康熙见过太子后心情都好了，“胸中亦不更有郁结矣”。但这里提到在召见废太子前，还召见了皇八子。皇八子刚因张明德案被削去了贝勒爵位，为什么又要在召见废太子前这么敏感的时间，先召见他？这也是在《实录》中无解的问题。


% 原书第57–58页

在这次会见前一周，《实录》在11月26日（四十七年十月十五日丁巳）这天记录了皇三子向康熙禀报，大皇子找了喇嘛“咒诅废皇太子”，并发现用来实施镇魇法术的物件。镇魇事件是康熙在短时间内决定复立太子的关键事件，这是以前学者都注意到的。从孟森开始，后来学者都同意镇魇事件是整个从废到立的转折点，但同时也都认为，镇魇事件只是康熙要复立太子时找到的一个借口。用孟森的话说：“乘此疑团，遂认为被厌胜，以图一时之复位。”也就是说孟森认为康熙在废黜太子以后后悔了，因而借着镇魇这样的巫术事件，把太子之前的问题说成是被巫术所害，从而达到复立太子的目的。\NoteRef{02}{48}


% 原书第58页

是什么原因让现当代学者认为镇魇事件只是康熙为复立太子找的借口？大概有二。一是传统史料没有关于镇魇事件的详细记载，因而只能对该事件进行推断。二是现当代学者由于有科学学习的背景，对“镇魇”这种巫术相关的记载历来都抱持高度怀疑的态度。


% 原书第58–59页

纪理安对镇魇事件的全过程做了详细记录。他认为镇魇事件就是康熙真心想要复立太子的转折点，并不是借口。当康熙得知镇魇后，他的第一反应，总结起来类似于我们今天知道投毒案后的反应，首先关注的是如何破除镇魇，也就是如何解毒。纪理安的记录中没有提到过皇三子举报的事，而说康熙最早是从大皇子身边一个杨姓太监的口供中得知：大皇子与喇嘛有来往，有镇魇之事。康熙听到这个报告首先想到的不是要处理大皇子，而是要从他那里找到“解药”。当时人知道，对于“造魇魅符书咒诅”等巫术，必须弄清施法细节后才有可能破解。这和投毒案中，必须第一时间知道毒药细节才能解毒是一样的。纪理安也领会到了找出镇魇施法的细节是整个案件的难点，因而他在记录中特别赞赏了康熙对破案的布置。\NoteRef{02}{49}


% 原书第59页

康熙为了避免打草惊蛇，首先安排了几队人马，同时抓捕了与案件相关的四名喇嘛。然后，他指示在傍晚之时把四个喇嘛带到大皇子的府中。纪理安认为康熙这一出其不意的安排起到了关键作用。大皇子和喇嘛在这样突然会面的压力下，立马都招了供，并交代了所有的施法细节，其中包括施法用的人偶所埋的具体位置。审问的人在得知位置后，连夜带人挖掘，分别在京城几处地方把人偶从地下挖了出来，并立即秘密送到康熙手中。\NoteRef{02}{50}

% 读者评论：02 第二章 一废太子.md，第5条
\begin{comment}
划重点：

康熙听到这个报告首先想到的不是要处理大皇子，而是要从他那里找到“解药”。

他（康熙）指示在傍晚之时把四个喇嘛带到大皇子的府中。

分别在京城几处地方（提示：不一定是太子、八爷府，有可能是法术阵眼）把人偶从地下挖了出来，并立即秘密送到康熙手中。
\end{comment}


% 原书第59–60页

纪理安还记录了康熙见到这些人偶后的反应。第一个人偶从外形上一看就是针对康熙的。这让康熙明白大皇子密谋的施法对象还包括自己，这完全出乎了他的意料。他之前只猜到大皇子针对太子，万万没有想到连他这位父亲也不放过。康熙之前知道大皇子有这样那样的问题，但是总体上认为与大皇子之间的父子感情是紧密的，他对大皇子十分信任：他不仅很早就给大皇子封了郡王，还把内务府很多事务交给他办理。几乎每一次外出，康熙都会带上大皇子。他们还曾一起远征大漠，是名副其实的“上阵父子兵”。所以当康熙得知大皇子施巫术的对象还包括自己时，既惊讶又愤怒。而对于镇魇的目标还包括康熙这一点，《实录》并没有提及。


% 原书第60页

纪理安还分析说，针对康熙的人偶让康熙相信了他是因为被施法镇魇后，才会“憎恨太子”，“自己才会失去理智”\NoteRef{02}{51}同时这一段时间，康熙的身体也确实不好。纪理安称康熙失眠，以致身体都有一些浮肿。\NoteRef{02}{52}人偶也让康熙联想到自己这段时间生病的根源。


% 原书第60页

康熙在气愤后，也有释然，因为人偶都已经挖了出来，算是找到了病因。因而康熙立马安排了下一步：安排喇嘛做法事来破解。两天以后，即12月5日，喇嘛入内廷做法，解除镇魇。纪理安说当天上午，他在内廷和一个管理钟表的太监在一起，法器演奏和念经的声音都听得一清二楚。\NoteRef{02}{53}


% 原书第60页

纪理安还提到一个重要的细节，从挖出来的人偶可知，除了康熙和太子以外，大皇子镇魇的目标还有皇八子。这也就解释了为什么前一天（12月4日），康熙召见太子前要专门召见皇八子。因为第二天，将会有喇嘛来施法为他们父子三人破除镇魇。当天，喇嘛把镇魇中的邪魔部分从父子三人身上转移到了事先找来的三名囚犯身上。法术完毕后，宫里派人把这三名囚犯送出直隶地界，并命令他们永不准回京。\NoteRef{02}{54}


% 原书第60–61页

在纪理安的文字中，康熙不是借镇魇来复立太子，而是真正参与到了镇魇的破解中，相信镇魇使自己以及太子做出了非同寻常之事。镇魇案以及镇魇后康熙与太子含泪相见的场景在当时已被社会所周知，所谓“国人所有闻之者，无不心碎落泪”\NoteRef{02}{55}。最起码，纪理安相信镇魇以及前后的因果关系，这也是为什么在他的叙述中镇魇是转折点。正因为康熙深信镇魇之术，才会对大皇子做出最重的惩罚：完全收回了在物质和名誉上给予大皇子的封赏，不仅废了他的王爵，还把分派给他管理的上三旗的牛录（八旗中的编制）全部收回。康熙还指示下属要紧盯大皇子，以免他再次“施法诅咒”。\NoteRef{02}{56}对大皇子的处罚，纪理安的记录与《实录》中的记载基本一致。两种记录都提到，从大皇子那里收缴来的财产，一部分给了大皇子的一个儿子，一部分给了皇十四子。


\begin{sectionnotes}

% 原书注释：原书第396页
\BookNote{02}{47}{《圣祖仁皇帝实录》，卷235。}

% 原书注释：原书第396页
\BookNote{02}{48}{孟森，《明清史讲义》，463。}

% 原书注释：原书第396页
\BookNote{02}{49}{AP3, 511-513.}

% 原书注释：原书第396页
\BookNote{02}{50}{AP3, 512.}

% 原书注释：原书第396页
\BookNote{02}{51}{AP3, 513.}

% 原书注释：原书第396页
\BookNote{02}{52}{AP3, 486.}

% 原书注释：原书第396页
\BookNote{02}{53}{AP3,518-519.}

% 原书注释：原书第396页
\BookNote{02}{54}{AP3,518-519.}

% 原书注释：原书第396页
\BookNote{02}{55}{《康熙朝满文朱批奏折全译》，1651。}

% 原书注释：原书第396页
\BookNote{02}{56}{《康熙朝满文朱批奏折全译》，1651。}

\end{sectionnotes}
